% !TeX root = ../教材题库_全部条目_转写工作稿.tex
% Source: XJ-B1-C06-LAB-P234-Y0145
% 第6章 / 6.3 统计图表 / 教材 p.234 / PDF 第241页
% OCR未校对：公式与文字均可能有误；未完成数学TeX转写；禁止直接发给学生
% 原题及相邻原解截图（核对用）：../题目原图/XJ-B1-C06-LAB-P234-Y0145-01.webp
% 本文件只在转写完成后作为题干使用；OCR 条目尚不是可靠题干。
\begin{bankitem}{利用计算机制作统计图表}{第6章 / 6.3 统计图表 / 教材 p.234 / PDF 第241页}{XJ-B1-C06-LAB-P234-Y0145}{OCR未校对：公式与文字均可能有误；未完成数学TeX转写；禁止直接发给学生}
\begin{ocrdraft}
[教材p.234 / PDF第241页]
\par
数学实验
\par
利用计算机制作统计图表
\par
当需要处理大量数据或绘制准确度较高的统计图表时，利用计算机软件是一种
\par
方便又快捷的方法．ExceI是其中一种最常用的绘制统计图表的软件，可以用来做
\par
数据统计、制表、数据分析等．
\par
我们以6．3节案例1的样本数据为例，利用Excel生成频数分布表．首先将表
\par
6一1的表头和数据，以行（或列）的方式整体输入ExceI中，然后用鼠标选中这些数
\par
据单元格，单击〖数据〗菜单中的〖数据透视表和数据透视图〗，根据对话框的提示，
\par
单击完成就可得到“数据透视表”，如图1为不同类型饮料和顾客性别的频数分布表．
\par
B
\par
C
\par
F
\par
G
\par
2
\par
酸
\par
计数项：饮料类型顾客性别@
\par
3
\par
女
\par
茶饮料
\par
饮料类型@
\par
男
\par
女
\par
总计
\par
4
\par
矿泉水
\par
茶饮料
\par
4
\par
7
\par
5
\par
女
\par
矿泉水
\par
果汁
\par
2
\par
4
\par
6
\par
6
\par
女
\par
碳酸饮料
\par
矿泉水
\par
6
\par
4
\par
7
\par
男
\par
矿泉水
\par
其他
\par
3
\par
5
\par
8
\par
酸饮料
\par
碳酸饮料
\par
9
\par
6
\par
15
\par
9
\par
女
\par
茶饮料
\par
总计
\par
24
\par
26
\par
10
\par
女
\par
果汁
\par
凵
\par
男
\par
碳酸饮料
\par
12
\par
男
\par
碳酸饮料
\par
女
\par
果汁
\par
13
\par
女
\par
矿泉水
\par
14
\par
女
\par
矿泉水
\par
巧
\par
根据数据分析的需求，我们还
\par
茶饮料
\par
可以利用ExceI绘制不同类型的统
\par
男
\par
计图，比如条形图、折线图、扇形
\par
女
\par
图、散点图等，其操作步骤大致相
\par
碳酸饮料
\par
果汁
\par
0
\par
同：选中输入的数据，单击工具栏
\par
（或〖插入〗菜单）中的“图表”标
\par
识，按照图表向导进行相应操作，即
\par
可得到不同类型的统计图（如图2）．
\par
其他
\par
矿泉水
\par
请同学们借助软件绘制6．3节案例
\par
2、案例3的统计图．
\end{ocrdraft}
\end{bankitem}
